\subsection{A Worked Job-Service Protocol}

The complete solution is \filepath{tutorial/examples/03-jobs.abclr2}.
It uses two workers to show that completion cannot be inferred from
dispatch order. Each input is a two-element list \code{[id value]},
where \code{id} is an integer. Each worker sends exactly one
\code{[:done id status result]} reply, including when the value is invalid.
This rule lets the monitor count completions without confusing a rejected
job with a missing response.

\noindent\begin{minipage}{\linewidth}
\begin{tabular}{L{0.24\textwidth}L{0.40\textwidth}L{0.26\textwidth}}
\textbf{Receiver} & \textbf{Message} & \textbf{Reply destination}\\\hline
Dispatcher & \code{[:run jobs]} & The original caller.\\
Monitor factory & \code{[:new count destination]} & The dispatcher;
returns a fresh monitor object.\\
Worker & \code{[:square id value]} & The batch's monitor.\\
Monitor & \code{[:done id status result]} & After the last completion,
sends a report to the original caller.\\
\end{tabular}
\end{minipage}

The workers deliberately use the same behavior; the second definition
creates another independent object:

\wholelisting[firstline=2,lastline=14]{tutorial/examples/03-jobs.abclr2}

The monitor factory captures the batch's final destination.
Each created monitor has its own counters and result list:

\wholelisting[firstline=16,lastline=33]{tutorial/examples/03-jobs.abclr2}

The dispatcher creates the monitor before submitting any work.
It routes even job identifiers to one worker and odd identifiers
to the other, keeping \code{@ monitor} on every request:

\sourcelisting[firstline=35]{tutorial/examples/03-jobs.abclr2}

The caller waits on the dispatcher, but the dispatcher handler ends after
dispatching. The monitor sends the eventual reply to the caller's \code{R}.
The empty case replies immediately because there will be no worker
completions to trigger the monitor.

\begin{codeblock}
\begin{lstlisting}
(inc "tutorial/examples/03-jobs" :use-source t)
[dispatcher <== [:run [[1 2] [2 3] [3 :bad]]]]
; => (:REPORT 2 1 ((... completion entries ...)))
[dispatcher <== [:run nil]]
; => (:REPORT 0 0 NIL)
\end{lstlisting}
\end{codeblock}

The completion entries are \code{(1 :OK 4)}, \code{(2 :OK 9)}, and
\code{(3 :FAILED :NOT-A-NUMBER)}. Their order is not part of the interface.
The automated check sorts entries by identifier before comparing them.
The counts, identifiers, statuses, and computed values are what the
protocol promises.

\subsection{Invariants and Limits}

For every monitor, \code{remaining + successful + failed} equals the
number of submitted jobs. A reply is sent only when \code{remaining}
reaches zero. Because each worker responds exactly once in this example,
the final report means the whole batch has completed.

The demonstration assumes well-formed pairs with integer identifiers,
no duplicate completion messages, and no worker failures outside the
declared invalid-value branch. Production extensions would need
duplicate detection, timeout or cancellation policy, and validation
of malformed submissions. A timeout is a policy decision, not proof
that a remote computation stopped.

\section{Worked Exercise Solutions}

Run \filepath{07-solutions.abclr2} after attempting the exercises:

\begin{lstlisting}
(inc "tutorial/examples/07-solutions" :use-source t)
\end{lstlisting}

\subsection{Exercise 1: Sum of Squares}
\label{sol:squares}

\sourcelisting[firstline=3,lastline=9]{tutorial/examples/07-solutions.abclr2}

\begin{lstlisting}
[sum-of-squares <== [:sum [2 3 4]]]   ; => 29
[sum-of-squares <== [:sum nil]]       ; => 0
\end{lstlisting}

The temporary accumulator belongs to one invocation.
An empty \code{dolist} leaves it at its initial value, so the result
is zero without a special case.

\paragraph{Back to the original.}
In Section~\ref{sec:original-future}, \code{tree2leaves} also keeps
handler-local working data, but rewrites the remaining tree instead of
adding to an accumulator. Why is a fixed \code{cdr} step not enough there?
\emph{Answer:} a matched clause may expose a nested head and tail as new
work. Follow the rewritten tree, not just the original list's next element.

\subsection{Exercise 2: An Account Factory and a Transfer}
\label{sol:transfer}

The factory returns an anonymous object for each \code{:new} request.
Its initializer is copied into a separate state binding.
Initial balances must be nonnegative integers; this is a precondition,
not a check performed by the factory. Deposit and withdrawal amounts
are checked by the account handlers.

\wholelisting[firstline=3,lastline=27]{tutorial/examples/02-accounts.abclr2}

\begin{lstlisting}
(inc "tutorial/examples/02-accounts" :use-source t)
[alice <== [:transfer bob 30]]     ; => :TRANSFERRED
[alice <== :balance]               ; => 70
[bob <== :balance]                 ; => 50
[alice <== [:transfer alice 10]]   ; => :REJECTED
\end{lstlisting}

The self-transfer guard avoids a synchronous wait on the current object.
The demonstration uses trusted recipients and a sequential caller.
The recipient is credited before the sender is debited, and the two
changes are not atomic together. Opposing simultaneous transfers can
form a wait cycle. A production transfer protocol would need reservations
and a defined commit/cancel procedure; changing one send to past-type
alone does not solve that problem.

\paragraph{Back to the original.}
Why can the inherited feather query in Section~\ref{sec:original-inherit}
forward a destination, while this transfer waits for the deposit result?
\emph{Answer:} the inherited script can answer the original caller itself;
the transfer handler still needs the recipient's result before deciding
whether to debit its own state and report success.

\subsection{Exercise 3: Count Only After Acknowledgment}
\label{sol:count}

The solution creates a fresh consumer for each run. Every item returns
its post-increment count to a future, and the producer consumes exactly
as many acknowledgments as it submitted items:

\wholelisting[firstline=11,lastline=25]{tutorial/examples/07-solutions.abclr2}

\begin{lstlisting}
[counted-producer <== [:run [2 3 4]]]   ; => 3
[counted-producer <== [:run nil]]       ; => 0
\end{lstlisting}

The final count is queried only after all increments have been acknowledged.
The local consumer prevents another batch from affecting the result.
Each \code{next-value} may wait, but no worker waits back on the producer,
so this protocol does not introduce a circular dependency.

\paragraph{Back to the original.}
Why does this client count replies, while \code{print-leaves} waits for
\code{:finish}? \emph{Answer:} this protocol submits a known number of
items and requires exactly one reply per item. The original tree walker
sends a stream of leaves whose length the consumer does not know in advance.
Pcontrol's two-child join (Section~\ref{sec:pcontrol-fib}) is another
known-count protocol; it must collect two answers, not merely the first.

\subsection{Exercise 4: A Tagged Stream}
\label{sol:stream}

\wholelisting[firstline=27]{tutorial/examples/07-solutions.abclr2}

\begin{lstlisting}
[square-stream-client <== [:collect 1 4]]
; => (1 4 9 16)
\end{lstlisting}

Values arrive inside \code{[:value x]}, while \code{:done} is a distinct
terminal message. The client unwraps only the value messages.
This separation remains useful even if the data domain later includes
keywords or lists.

\paragraph{Back to the original.}
Would renaming \code{:finish} to another keyword in the original tree
walker solve its marker collision? \emph{Answer:} no; that keyword could
also occur as a leaf. The value wrapper distinguishes the \emph{role}
of a message instead of reserving one possible data value.

\subsection{Exercise 5: Unless at the Meta-Level}
\label{sol:unless}

Add this handler before the choice evaluator's fallback:

\sourcelisting[firstline=19,lastline=24]{tutorial/examples/05-reflection.abclr2}

When the predicate is true, reply to \code{C} with \code{nil}.
Otherwise evaluate the compiled body in its captured environment.
The test object's \code{:unless} handler invokes this protocol:

\begin{codeblock}
\begin{lstlisting}
[choice-test <== [:unless 3]]   ; => 3
[choice-test <== :selected]     ; => :UNLESS
[choice-test <== [:unless 8]]   ; => NIL
[choice-test <== :selected]     ; => :UNLESS (unchanged)
\end{lstlisting}
\end{codeblock}

The unchanged state verifies that the body was skipped for an even input.
Both control-flow paths still send a result to the caller's continuation.

\paragraph{Back to the original.}
Compare \code{my-if} (Section~\ref{sec:original-if}) and
\code{with-interrupt} (Section~\ref{sec:interrupt}). Where does each
reified request go? \emph{Answer:} \code{my-if}, like our choice form,
uses the evaluator; \code{with-interrupt} uses the application's metaobject.
For either route, trace the saved continuation separately from the
application handler's eventual \code{!} reply.
